\documentclass[10pt]{article}
\usepackage[T1]{fontenc}
\usepackage[utf8]{inputenc}
\usepackage[french]{babel}
\usepackage{amsmath,amssymb,mathrsfs,mathpazo}
\usepackage{geometry}
\geometry{a4paper,margin=2cm}
\usepackage{graphicx,array,tabularx,tabularray,colortbl,multirow,multicol,booktabs,enumitem,pifont,ulem,fancybox,diagbox,cancel}
\usepackage{pstricks,pstricks-add,pst-plot,pst-tree,pst-node,pst-text,pst-func,pst-math,pst-bspline,pst-3dplot,pst-eucl,pst-solides3d}
\let\pscalculate\undefined
\usepackage{pst-bezier}
\usepackage{tikz,pgfplots,tkz-tab}
\usetikzlibrary{arrows,arrows.meta,patterns,calc,positioning,trees,decorations.markings,shapes.geometric,shapes.misc,shapes.arrows,intersections,decorations.pathmorphing,decorations.pathreplacing,backgrounds,shadows,babel}
\pgfplotsset{compat=1.18}
\usepgfplotslibrary{fillbetween}
\usepackage[np]{numprint}
\usepackage{listings,makecell,pas-tableur,twoopt}
\let\barre\undefined
\usepackage[pstricks]{tablvar}
\definecolor{darkgreen}{RGB}{0,100,0}
\providecommand{\red}{\color{red}}
\providecommand{\blue}{\color{blue}}
\providecommand{\green}{\color{green}}
\providecommand{\black}{\color{black}}
\providecommand{\white}{\color{white}}
\providecommand{\gray}{\color{gray}}
\providecommand{\cyan}{\color{cyan}}
\providecommand{\magenta}{\color{magenta}}
\providecommand{\yellow}{\color{yellow}}
\providecommand{\violet}{\color{violet}}
\providecommand{\orange}{\color{orange}}
\providecommand{\euro}{\texteuro}
\providecommand{\EUR}{\texteuro}
\providecommand{\ieme}{\textsuperscript{e}}
\providecommand{\degres}{\ensuremath{^\circ}}
\providecommand{\degree}{\ensuremath{^\circ}}
\providecommand{\vv}[1]{\overrightarrow{#1}}
\providecommand{\Coord}[1]{#1}
\providecommand{\decosix}{}
\providecommand{\starredbullet}{\ensuremath{\bullet}}
\pagestyle{empty}
\setlength{\parindent}{0pt}
\usepackage{adjustbox}\DeclareUnicodeCharacter{2236}{:}\DeclareUnicodeCharacter{2190}{\ensuremath{\leftarrow}}
\providecommand{\R}{\mathbb{R}}
\providecommand{\N}{\mathbb{N}}
\providecommand{\Z}{\mathbb{Z}}
\providecommand{\Q}{\mathbb{Q}}
\providecommand{\D}{\mathbb{D}}
\providecommand{\C}{\mathbb{C}}
\providecommand{\e}{{\mathrm{e}}}
\providecommand{\vect}[1]{\overrightarrow{#1}}
\providecommand{\vectt}[1]{\overrightarrow{\mathrm{#1}}}
\providecommand{\vvt}[1]{\overrightarrow{\mathrm{#1}}}
\providecommand{\np}[2][]{{#2~#1}}
\providecommand{\vv}[1]{{\overrightarrow{#1}}}
\renewcommand{\d}{{\mathrm{d}}}
\providecommand{\degres}{{^\circ}}
\providecommand{\degre}{{^\circ}}
\providecommand{\blue}{\color{blue}}
\providecommand{\red}{\color{red}}
\providecommand{\enskip}{\quad}
\providecommand{\llbracket}{\lbrack\!\lbrack}
\providecommand{\rrbracket}{\rbrack\!\rbrack}
\providecommand{\euro}{{\text{€}}}
\providecommand{\no}{{\text{n°}}}
\providecommand{\interval}[2]{{[#1;#2]}}
\providecommand{\intervFO}[2]{{[#1;#2[}}
\providecommand{\dotfill}{\ldots}

\providecommand{\ds}{\displaystyle}
\providecommand{\textding}[1]{#1}
\providecommand{\Oij}{\(\mathrm{O};\vec{\imath},\vec{\jmath}\)}
\providecommand{\Oijk}{\(\mathrm{O};\vec{\imath},\vec{\jmath},\vec{k}\)}
\providecommand{\Ouv}{\(\mathrm{O};\vec{u},\vec{v}\)}
\providecommand{\intervFF}[2]{\ensuremath{[#1;#2]}}
\renewcommand{\euro}{\ifmmode\text{\texteuro}\else\texteuro\fi}
\DeclareUnicodeCharacter{00B0}{\ifmmode^{\circ}\else\textdegree\fi}
\DeclareUnicodeCharacter{20AC}{\ifmmode\text{\texteuro}\else\texteuro\fi}
% Énoncés : Éducation nationale. Transcription et corrigés : APMEP, Denis Vergès et contributeurs. Découpage et classement : CoopMaths, Sébastien Lozano, Éric Elter et contributeurs. Adaptation de la présentation : Maths Explorer.
% CC BY-SA : https://coopmaths.fr/www/about/
% https://coopmaths.fr/alea/static/bac/2023/tex/bac_2023_08_sujet1_caledonie_4_cor.tex
\begin{document}
\noindent\textit{APMEP / CoopMaths. CC BY-SA. Présentation adaptée par Maths Explorer.}\par\medskip
\begin{enumerate}
	\item
	\begin{enumerate}
		\item $\begin{gathered}\lim\limits_{\substack{x \to 0\\x>0}}f(x) = \lim\limits_{\substack{x \to 0\\x>0}}\left[5x^2 + 2x - 2x^2\ln(x)\right]\end{gathered}$
		
		\[\left.
		\begin{array}{r}
			\lim\limits_{x \to 0}\left(5x^2+2x\right) = 0\\
			\text{Par croissance comparée}\: \lim\limits_{\substack{x \to 0\\x>0}}x^2\ln(x) = 0
		\end{array}\right\} \overset{\substack{\text{par} \\ \text{somme}}}{\implies} \lim_{\substack{x \to 0\\x>0}}\left[5x^2 + 2x - 2x^2 \ln (x)\right] = 0\]
		
		\[\boxed{\lim_{\substack{x \to 0\\x>0}}f(x) = 0}\]
		\item $\lim\limits_{x \to +\infty}f(x) = \lim\limits_{x \to +\infty}\left[5x^2+2x-2x^2\ln (x)\right] = \lim\limits_{x \to +\infty}x^2\left[5+\dfrac{2}{x}-2\ln(x)\right]$
		
		\[\renewcommand{\arraystretch}{2}
		\left.
		\begin{array}{r}
			\lim\limits_{x \to +\infty}\dfrac{2}{x} = 0 \implies \lim\limits_{x \to +\infty}\left(5+\dfrac{2}{x}\right) = 5 \\
			\lim\limits_{x \to +\infty}\ln(x) = +\infty \implies \lim\limits_{x \to +\infty}-2\ln(x) = -\infty
		\end{array}\right\} \overset{\substack{\text{par} \\ \text{somme}}}{\implies} \lim_{x \to +\infty}\left[5 + \dfrac{2}{x} - 2\ln (x)\right] = -\infty
		\]
		
		\[\renewcommand{\arraystretch}{2}
		\left.
		\begin{array}{r}
			\lim\limits_{x \to +\infty}x^2 = +\infty \\
			\lim\limits_{x \to +\infty}\left[5 + \dfrac{2}{x} - 2\ln (x)\right] = -\infty
		\end{array}\right\} \overset{\substack{\text{par} \\ \text{produit}}}{\implies} \lim_{x \to +\infty}x^2\left[5+\dfrac{2}{x}-2\ln(x)\right] = -\infty
		\]
		
		\[\boxed{\lim_{x \to +\infty}f(x) = -\infty}\]
	\end{enumerate}
	\item\
	\begin{eqnarray*}
		f'(x) &=& 5 \times 2x + 2 - 2 \left[ 2x \times \ln(x) + x^2 \times \dfrac{1}{x}\right] = 10x + 2 - 2\left[2x\ln(x) + x\right] = 10x + 2 - 4x\ln(x) - 2x\\
		&=& \boxed{8x + 2 - 4x\ln(x)}
	\end{eqnarray*}
	\item
	\begin{enumerate}
		\item
		\[f''(x) = 8 - 4\left[1 \times \ln(x) + x \times \dfrac{1}{x}\right] = 8 - 4\ln (x) - 4 = 4 - 4\ln(x) = \boxed{4\left[1-\ln(x)\right]}\]
		\item La courbe $\mathcal{C}_f$ est au-dessus de ses tangentes si, et seulement si, $f$ est convexe, ce qui équivaut à $f''(x)$ est positif.
		
		$f''(x) \geqslant 0 \iff 4\left[1-\ln(x)\right] \geqslant 0 \iff 1-\ln(x) \geqslant 0 \iff -\ln(x) \geqslant -1 \iff \ln(x) \leqslant 1 \iff x \leqslant \text{e}$ (car la fonction ln est strictement croissante sur $]0~;~+\infty[$).\\
		\smallskip
		
		La courbe $\mathcal{C}_f$ est donc au-dessus de ses tangentes sur l'intervalle $]0~;~\text{e}]$.
		\item\
		\begin{center}
		\begin{tikzpicture}
		\tkzTabInit[lgt=2.5,espcl=3]{$x$/1,Signe de $f''(x)$/1,Variations de $f'$/2.5}{$0$,e,$+\infty$}
		\tkzTabLine{d,+,z,-,}
		\tkzTabVar{D-/$2$,+/$4\text{e}+2$,-/$-\infty$}
		\end{tikzpicture}
		
		$f'(\text{e}) = 8\text{e} + 2 - 4\text{e} \ln(\text{e}) = 8\text{e} + 2 - 4\text{e} = 4\text{e} + 2$
		\end{center}
	\end{enumerate}
	\item
	\begin{enumerate}
		\item Sur $]0~;~\text{e}]$, la fonction $f'$ est strictement croissante avec $\begin{gathered}\lim\limits_{\substack{x \to 0\\x>0}}f'(x) = 2\end{gathered}$, donc pour tout $x \in ]0~;~\text{e}]$,on a $f'(x) > 0$.\\
		\smallskip
		
		Sur $[\text{e}~;~+\infty[$, la fonction $f'$ est continue (puisque dérivable) et strictement décroissante. De plus, $f(_text{e}) = 4\text{e} + 2 > 0$ et $\lim\limits_{x \to +\infty}f'(x) = -\infty$. $0 \in ]-\infty~;~4\text{e}+2[$, donc, d'après le corollaire du théorème des valeurs intermédiaires, l'équation $f'(x) = 0$ admet une solution unique $\alpha$ sur $[\text{e}~;~+\infty[$.\\
		\smallskip
		
		Au final, l'équation $f'(x) = 0$ admet une solution unique $\alpha$ sur $]0~;~+\infty[$.\\
		\medskip
		
		On a :
		\[\boxed{\np{7.87} < \alpha < \np{7.88}}\]
		\item On sait que sur $]0~;~\text{e}]$ on a $f'(x) > 0$. Puisque $f'$ est décroissante sur $[\text{e}~;~+\infty[$ et qu'elle s'annule en $\alpha$, $f'(x) \geqslant 0$ sur $[\text{e}~;~\alpha]$ et $f'(x) \leqslant 0$ sur $[\alpha~;~+\infty[$. Donc $f'(x)$ est positif sur $]0~;~\alpha]$ et négatif sur $[\alpha~;~+\infty[$.
		
		\begin{center}
		\begin{tikzpicture}
		\tkzTabInit[lgt=2.5,espcl=3]{$x$/1,Signe de $f'(x)$/1,Variations de $f$/2.5}{$0$,$\alpha$,$+\infty$}
		\tkzTabLine{d,+,z,-,}
		\tkzTabVar{D-/$0$,+/$f(\alpha)$,-/$-\infty$}
		\end{tikzpicture}
		\end{center}
	\end{enumerate}
	\item On a $f''(\alpha) = 8\alpha + 2 - 4\alpha \ln(\alpha)$, donc :
	\begin{eqnarray*}
		f'(\alpha) = 0 &\iff& 8\alpha + 2 - 4 \alpha \ln(\alpha) = 0\\
		&\iff& 4\alpha\ln(\alpha) = 8\alpha + 2\\
		&\iff& \ln(\alpha) = \dfrac{2(4\alpha+1)}{4\alpha} \qquad (\alpha \neq 0)\\
		&\iff& \boxed{\ln(\alpha) = \dfrac{4\alpha+1}{2\alpha}}
	\end{eqnarray*}
	
	On a donc :
	\begin{eqnarray*}
		f(\alpha) &=& 5\alpha^2 + 2\alpha - 2\alpha^2 \ln(\alpha) \\
		&=& 5\alpha^2 + 2\alpha - 2\alpha^2 \times \dfrac{4\alpha+1}{2\alpha} \\
		&=& 5\alpha^2 + 2\alpha - \alpha(4\alpha + 1) \\
		&=& 5\alpha^2 + 2\alpha - 4\alpha^2 - \alpha \\
		&=& \boxed{\alpha^2  + \alpha}
	\end{eqnarray*}
	\item D'après la question \textbf{4.a.}, on a :
	\begin{eqnarray*}
		\np{7.87} < \alpha < \np{7.88} &\iff& \np{61.9369} < \alpha^2 < \np{62.0944} \\
		&\iff& \np{69.8069} < \alpha^2 + \alpha < \np{69.9744} \\
		&\iff& \boxed{\np{69.8069} < f(\alpha) < \np{69.9744}}
	\end{eqnarray*}
\end{enumerate}

\end{document}
